% Vérifié (19 oct. 2026) : citations et renvois contrôlés sur lecture/p1 (147-186) ; « souvenirs affluent » et « 2 fenêtres » (f. 83 v°) rétablis ; « lgt » (lecture incertaine) paraphrasé ; « trois jours désœuvrés » corrigé ; consigne du 8 mars datée ; Wentworth et Monin ramenés à l'hypothèse ; dernière mention de M. Calmels datée (f. 71 v°) ; pages OH (Bapaume, Arras) alignées sur R4 ; paragraphe sur la crise de mars allégé (Levy et 21 avril déjà dans intro_interprete), note Blenkinsop non vérifiable supprimée, note SHD reformulée ; fin allégée (destruction des papiers déjà dans intro_biographie ; « nouvelle permission » non sourcée retirée).
% chapitre_VII.tex -- introduction du chapitre VII (clichés 147-186 ; ff. 73 v°-93 r°), fragment LaTeX sans préambule.
% Sources : carnet (lecture/p1/147-186), commentaires historiques (commentaire/147-151, 176-181), dossiers R4, R2, R3, R7, R8.
% Rédaction du 18 oct. 2026.

Le 7~février 1917 au soir, de retour de permission, Hippolyte Combes s'installe à Hérissart, village du plateau picard à trois kilomètres de Toutencourt, où siège depuis juillet 1916 le quartier général de la Fifth Army du général Gough\footnote{C.~Baker, \textit{The Long, Long Trail}, «~The British Fifth or Reserve Army 1916--1918~» (copie Wayback du 23~juillet 2024) ; distance calculée d'après les coordonnées des mairies.}. Le capitaine Williams y a transporté en son absence le bureau du \textit{Purchase Board}. Le travail reste le même ; son cadre change. L'interprète relève désormais du bureau de la Mission française de l'armée, et non plus de celui d'un corps ; le front, que la 18th Division tenait au sud de l'Ancre, s'éloigne vers l'est avec le repli allemand de mars ; l'arrière-front se vide de ses troupes. Son poste lui-même manque de disparaître : le 10~mars, le départ du capitaine laisse Combes «~sans fonctions~», et il ne doit qu'à l'intervention du major Levy d'être maintenu. Le nom de Marguerite Calmels, la fiancée, n'apparaît plus après le retour (dernière mention au f.~71~v\textsuperscript{o}). Le carnet tient le relevé des courses, des ciels et des paysages d'un pays que la guerre quitte.

\subsection*{Le carnet}

Le Town Major attribue à Combes une «~petite chamb, no~7b~», au premier étage d'un café d'angle, avec vue sur la place et l'église (f.~73~v\textsuperscript{o}) ; bureau et mess occupent le n\textsuperscript{o}~104, «~mauvaise installation~» (f.~74~r\textsuperscript{o}). Le lendemain, par un grand gel, il marche jusqu'à Rubempré et trouve la région «~moins banale qu'environs d'Acheux~», avec ses «~petites hauteurs, couronnées de moulins à vent~» (f.~74~v\textsuperscript{o}). Le 9, au bureau de la Mission à Toutencourt, il apprend qu'il sera «~administré par l'armée et payé les 3 et 18 de chaque mois~» (f.~75~r\textsuperscript{o}). Le 10, une tournée à Martinsart le conduit, route d'Aveluy, auprès de son «~ancienne Cy~», la 152nd Company, «~campée champ devasté sous des tentes~» ; il déjeune au mess, cause avec les sous-officiers et les hommes, puis passe au bureau de la Mission du II Corps, à Senlis, où il juge le capitaine de Casteljau «~énorme \& bête~» (f.~75~v\textsuperscript{o}). Le 12, au retour de Méaulte, il traverse Ribemont et Méricourt, «~ces villages où vecu heures bien vide \& triste~» (f.~76~v\textsuperscript{o}). Suivent les tournées «~de ferme en ferme pour trouver paille \& foin~», à Authie, Beaucourt, Fréchencourt, Bavelincourt et Pierregot (ff.~77~r\textsuperscript{o}--78~v\textsuperscript{o}), par des routes que le dégel change en «~fondrières~» (f.~78~r\textsuperscript{o}). Le 19, il passe dans l'ancienne chambre du capitaine, dont les briques «~ruissellent d'eau~» (f.~78~v\textsuperscript{o}). L'humeur suit le temps : «~me suis tt à fait ressaisi~» le 18 (f.~78~r\textsuperscript{o}), «~Impressions de solitude \& detresse morale~» le 21 (f.~79~r\textsuperscript{o}), puis des journées presque désœuvrées dans un «~Village où ns sommes exilés, absolt perdu sur hauteurs~» (f.~79~v\textsuperscript{o}).

Mars s'ouvre sur l'incertitude. Le 5, Combes apprend l'arrivée prochaine d'un officier australien et d'un interprète «~pour ns remplacer~» et se dit «~Très content de quitter le Cap~» (f.~81~r\textsuperscript{o}). L'officier, «~plutôt vulgaire~», et le nouvel interprète arrivent le 6 ; au bureau de la Mission, un maréchal des logis lui promet «~bon poste~» (f.~81~v\textsuperscript{o}) ; le 8, on lui prescrit de rester auprès du capitaine jusqu'à son départ et de se rendre ensuite au dépôt de Toutencourt (f.~82~r\textsuperscript{o}). Williams part le 10 au matin. Combes, «~sans fonctions, archi-libre~», note : «~souvenirs affluent de ce 6 derniers mois~» et s'inquiète d'une «~situation precaire~» (f.~82~v\textsuperscript{o}). Le 11, ses sacs déjà déposés au mess des interprètes de Toutencourt, il apprend que, «~sur demande instante du Major Levy~», il restera au \textit{Purchase Board} ; il s'installe le soir même au billet~220 d'Hérissart, chambre calme à deux fenêtres et «~2 Tables pour écrire~» (ff.~83~r\textsuperscript{o}--v\textsuperscript{o}). Le service reprend le 12 avec l'officier australien, peut-être le lieutenant Wentworth nommé à partir du 17 (f.~84~v\textsuperscript{o}), dans une Wolseley conduite par un Australien, «~moins confortable que le Vauxhall~» d'autrefois (f.~83~v\textsuperscript{o}) ; bureau et mess passent le 14 au billet~216, en face, «~1000 fois mieux~» (f.~84~r\textsuperscript{o}). La veille, Combes a rendu visite à «~M. \& Mme Flutre, instituteurs, au Mesnil~», dont le fils est «~eleve de Norm\textsuperscript{le} Sup\textsuperscript{re}~» (f.~84~r\textsuperscript{o}).

Les tournées s'étendent alors vers l'Ancre. Le 16, à Albert, Combes revoit les maisons touchées par les obus et la basilique, puis s'arrête longuement à Ribemont, où il retrouve des «~connaissances trè sympathiques~» (f.~84~v\textsuperscript{o}). Le 19, des officiers de la 63rd (Royal Naval) Division, revenus des tranchées de Miraumont, lui donnent des «~renseigts très intéressants~» (f.~85~r\textsuperscript{o}) ; le 21, entre Contay, Warloy et Senlis, les routes sont «~encombrées par troupes moving~» (f.~85~v\textsuperscript{o}). Le 26, Monin, sans doute un autre interprète, cantonne une nuit au village avec sa compagnie : Combes le trouve «~vieilli, enlaidi, usé~» (f.~86~v\textsuperscript{o}). Le 27, les troupes partent à l'aube (f.~86~v\textsuperscript{o}) ; le 29, «~Village parait abandonné. Plus de troupes. Rues désertes. Calme délicieux~», et le 31, sous la grêle : «~J'aime ce calme \& ce silence~» (f.~87~r\textsuperscript{o}).

Avril est le mois des neiges tardives et des derniers achats. Combes passe chez le notaire de Toutencourt, M\textsuperscript{e}~Bienaimé, le 3 et le jour de Pâques (ff.~87~v\textsuperscript{o}, 89~r\textsuperscript{o}). Le 6, un membre du bureau de la Mission, venu du quartier général, paie la solde à Toutencourt ; le soir, devant le poêle, Combes cause longuement de l'Australie avec Wentworth (f.~88~v\textsuperscript{o}). Le 13, il converse longuement avec le sous-lieutenant Gilbert, «~de Jersey, il est très instruit et poète~» (f.~89~v\textsuperscript{o}), qui l'accompagnera jusqu'à la fin du mois. Le 19, il retrouve Bresle, Laviéville et Hénencourt «~vides et déserts~» ; le 20, l'auto s'embourbe à la sortie de Buire et des Australiens la dégagent : «~Quels types!~» ; à Ribemont, les Colart lui font «~meilleur accueil~» (f.~90~v\textsuperscript{o}). Les soirées se passent chez une famille d'Hérissart, les Denis (ff.~89~r\textsuperscript{o}, 90~r\textsuperscript{o}, 91~v\textsuperscript{o}, 92~r\textsuperscript{o}). Le 21, après une nuit de «~violente canonnade --- sud d'Arras~», Combes note : «~Cherché vainement paille à acheter~» (f.~91~r\textsuperscript{o}). Le 28, un officier revenu de Bournemouth annonce que l'on quittera bientôt le village pour Cambrai : «~Regretterai bp les points de vue et panoramas circulaires des environs~» (f.~92~v\textsuperscript{o}). Le 29, deux ressorts cassés à Contay le ramènent à Hérissart à pied puis sur une voiture à eau (f.~92~v\textsuperscript{o}) ; le 30, dernière page du carnet, il fait ses courses «~sans pardessus~» et se couche tôt, fatigué (f.~93~r\textsuperscript{o}).

Les comptes sont ceux d'une vie réglée : mess (33~f pour février, «~depuis le 8~», f.~80~r\textsuperscript{o}), blanchissage, coiffeur, cartes postales, un mandat de 42~f à M\textsuperscript{me}~Spagnol, de Nice, sans doute l'expéditrice des fleurs reçues à Toulouse (f.~79~r\textsuperscript{o}), une ceinture de cuir de 24~f~50 achetée à Doullens (f.~89~v\textsuperscript{o}) ; le 3~mars, la Mission retient 35~fr. «~pour 7 jours de permission~» (f.~81~r\textsuperscript{o}).

\subsection*{Le contexte}

Les deux premiers mois du chapitre sont ceux du dégel et du repli allemand. Le gel, commencé vers le 14~janvier, cède le 16~février ; le lendemain, la 18th Division attaque à Boom Ravine, près de Miraumont, dans une «~boue profonde~», et y perd 54~officiers et 1\,135~hommes, sept jours après la visite de Combes à son ancienne compagnie\footnote{C.~Falls, \textit{Military Operations, France and Belgium, 1917}, vol.~I, \textit{The German Retreat to the Hindenburg Line and the Battles of Arras}, Londres, Macmillan, 1940 (ci-après \textit{OH 1917}, I), p.~66 et 75--82 ; G.~H.~F.~Nichols, \textit{The 18th Division in the Great War}, Édimbourg et Londres, Blackwood, 1922, p.~139--143.}. Le 24~février, des patrouilles de la division découvrent le premier repli ennemi. L'opération \textit{Alberich}, ordonnée le 4~février, porte en quatre jours de marche, du 16 au 19~mars, les armées allemandes sur la ligne Hindenburg, en laissant derrière elles villages détruits, arbres abattus et puits comblés ; Bapaume est occupé le 17, et l'hôtel de ville, piégé, y saute dans la nuit du 25, tuant surtout des Australiens et deux députés français\footnote{\textit{OH 1917}, I, p.~95--101, 110--115, 131--132, 138--139 et 148.}. Le carnet n'en dit rien. Il note en revanche ce que l'histoire officielle constate à l'arrière : après le dégel, «~même les voitures légères crevaient la mince croûte d'empierrement~», un cinquième des camions de l'armée tombe en panne, et les routes sont «~presque ruinées~» après la mi-mars\footnote{\textit{Ibid.}, p.~103 et 151.} ; l'auto du \textit{Board} est «~aux workshops~» les 17 et 23~mars et le 18~avril (ff.~84~v\textsuperscript{o}, 86~r\textsuperscript{o}, 90~v\textsuperscript{o}), et les routes sont «~creusées de sillons profonds par gros camions~» (f.~87~r\textsuperscript{o}).

La 18th Division, quant à elle, prend Irles le 10~mars, poursuit l'ennemi jusqu'à Saint-Léger et, relevée le 20, quitte la Somme. Le 22, le 7th Bedfordshire, de la brigade que ravitaillait la 152nd Company, gagne Contay et défile à Senlis devant le commandant du corps\footnote{Nichols, \textit{op.~cit.}, p.~157--163 ; Bedfordshire Archives Service, X550/8/1/1703 : \textit{War Diary}, 7th (Service) Battalion, Bedfordshire Regiment, mars 1917 ; la 53rd Brigade passe par Warloy et Rainneville (T.~M.~Banks et R.~A.~Chell, \textit{With the 10th Essex in France}, 2\textsuperscript{e}~éd., Londres, Gay \& Hancock, 1924, p.~186--187).} , le jour même où Combes repasse par Senlis (f.~85~v\textsuperscript{o}) ; les «~troupes moving~» qu'il voyait la veille entre Contay et Senlis appartenaient peut-être à ce mouvement. La division se repose en avril entre Aire et Hazebrouck ; le 27, ses transports remontent vers Arras «~under Capt. West Comdg 152nd Coy A.S.C.~»\footnote{Bedfordshire Archives Service, X550/8/1/1704 : \textit{War Diary} cité, avril 1917, ordre n\textsuperscript{o}~5 du 26~avril ; Nichols, \textit{op.~cit.}, p.~165--166.}. L'arrière-front, lui, se dépeuple. Le 15~février, la Fifth Army avait reçu le I Anzac Corps\footnote{\textit{OH 1917}, I, p.~74.}, ce qui explique sans doute les officiers, chauffeurs et soldats australiens du carnet ; mais le quartier général de l'armée quitte Toutencourt pour Albert le 28~mars, celui du V Corps quitte Acheux le 2~avril, et les hôpitaux d'évacuation de Puchevillers et de Contay suivent le front vers l'est\footnote{Baker, \textit{The Long, Long Trail}, notice citée, «~The British V Corps in France and Flanders~» (copie Wayback du 5~août 2020) et «~Locations of British Casualty Clearing Stations~» (copie Wayback du 13~janvier 2018).}. Le silence d'Hérissart à la fin de mars en est la conséquence, et Combes le dit lui-même de Toutencourt le 6~avril, «~si tranquille, depuis départ H.Q. et troupes~» (f.~88~v\textsuperscript{o}). Le 9~avril s'ouvre la bataille d'Arras, dont Combes, à près de quarante kilomètres, ne note la canonnade que dans la nuit du 20 au 21, deux jours avant la deuxième bataille de la Scarpe\footnote{\textit{OH 1917}, I, chap.~VIII--XV ; distance calculée d'après les coordonnées des mairies.} ; la déclaration de guerre des États-Unis, le 6~avril\footnote{W.~Miles, \textit{Military Operations, France and Belgium, 1916}, vol.~II, Londres, Macmillan, 1938, p.~531, n.~4.}, n'y laisse pas davantage de trace.

La crise de mars éclaire la situation de l'interprète. La Mission paie, convoque et affecte : le calendrier de paie annoncé le 9~février est tenu à quelques jours près (18~février, 3 et 18~mars, 6 et 21~avril), y compris après le départ du quartier général pour Albert (ff.~78~r\textsuperscript{o}, 81~r\textsuperscript{o}, 85~r\textsuperscript{o}, 88~v\textsuperscript{o}, 91~r\textsuperscript{o}). L'officier britannique peut cependant obtenir le maintien d'un interprète apprécié ; aucun texte n'a été retrouvé sur la réduction du nombre des interprètes que Combes redoutait depuis l'automne\footnote{Dossier R3 de l'édition, F16--F17 ; les cartons du Service historique de la Défense susceptibles d'en traiter (GR~16~N~509 ; 17~N~447--448) n'ont pas été dépouillés.}.

Parmi les personnes du chapitre, peu sont identifiées. Le fils des instituteurs Flutre est très probablement Louis-Fernand Flutre (1892--1978), né à Mesnil-Martinsart, ancien élève de l'École normale supérieure, futur professeur à la faculté des lettres de Lyon et dialectologue du parler picard de son village\footnote{Wikidata, Q559663 ; L.-F.~Flutre, \textit{Le Parler picard de Mesnil-Martinsart (Somme)}, Genève et Lille, 1955. Mesnil-Martinsart étant alors proche du front, la visite, faite «~Après tea~» un jour où Combes est «~Resté seul à Herissart~», a sans doute lieu à Hérissart, où les parents seraient réfugiés (hypothèse).}. Le lieutenant Wentworth pourrait être l'un des deux officiers de l'armée australienne de ce nom présents en France au printemps de 1917, Henry Coates Wentworth, du train de la 2nd Australian Division, ou Fitzwilliam Wentworth, éleveur du Queensland\footnote{UNSW Canberra, \textit{The AIF Project}, notices 318876, 318872 et 318873.}. Le capitaine Williams, le sous-lieutenant Gilbert, les Denis, les Colart, chez qui Combes avait cantonné à Ribemont pendant l'hiver 1915--1916, et Monin ne le sont pas\footnote{Dossier R7 de l'édition, F18, F26, F50 et F68 ; pour Gilbert, aucun des cinq officiers de l'ASC de ce nom de la série WO~339 des National Archives n'est relié à Jersey, et son arme n'est pas connue.}.

L'annonce du 28~avril, «~ns irons à Cambrai~», suppose une avance que la guerre ne donnera pas : la ligne Hindenburg passe alors à douze ou treize kilomètres à l'ouest de la ville, qui restera allemande jusqu'au 9~octobre 1918\footnote{\textit{Guides illustrés Michelin des champs de bataille. St-Quentin -- Cambrai}, Clermont-Ferrand, Michelin, 1921, p.~62--63 ; distance calculée d'après les mairies de Mœuvres et d'Havrincourt.}. Le carnet s'arrête deux jours plus tard, au recto du dernier feuillet, sur une journée de printemps et sans mention de départ ; la couverture porte les mêmes bornes, «~Du 25 juillet 1916 au 30 avril 1917~». Le carnet renvoie lui-même à un «~Journal~» tenu sur «~feuille speciale~» (f.~71~r\textsuperscript{o}), mais aucune suite n'a été retrouvée, et les mois suivants ne sont connus que par d'autres documents, jusqu'au mariage célébré à Millau le 25~octobre 1917.
